\documentclass[11pt]{article}
\usepackage[T1]{fontenc}
\usepackage{lmodern}
\usepackage{amsmath,amssymb,amsthm,mathtools}
\usepackage[margin=1in]{geometry}
\usepackage[hidelinks]{hyperref}
\usepackage{enumitem}

\newtheorem{definition}{Definition}[section]
\newtheorem{proposition}[definition]{Proposition}
\newtheorem{theorem}[definition]{Theorem}
\newtheorem{remark}[definition]{Remark}
\newtheorem{example}[definition]{Example}
\newcommand{\Req}{\operatorname{Req}}
\newcommand{\Sat}{\operatorname{Sat}}
\newcommand{\Def}{\operatorname{Def}}
\newcommand{\Adm}{\operatorname{Adm}}
\newcommand{\Real}{\operatorname{Real}}
\newcommand{\Cl}{\operatorname{Cl}}
\newcommand{\Hist}{\mathcal{H}}
\newcommand{\Obs}{\mathcal{E}}
\newcommand{\Laws}{\mathcal{J}}

\title{ConstructiveRT v0.4: Original-Property Residual Preservation, Propagation, and Joining\\
\large Conditional Formalization and Open Proof Obligations}
\author{Joseph Fredrickson\\Working mathematical research notes}
\date{November 12, 2026\\Revision 0.1}

\begin{document}
\maketitle
\begin{abstract}
This note formalizes the first two proposed Constructive Residual Preservation and
Propagation results (CRPP-1 and CRPP-2). The foundational goal is to determine
consequences from original property occurrences, witnessed evidence, and established
joining relationships, without allocating independent absence-property operations or
materializing hypothetical successful intermediate constructions. The results are
\emph{conditional}, applying only within complete, explicitly declared joining laws.
They do not establish a universal physical interaction grammar or a speed advantage
over conventional evaluation.
\end{abstract}

\section{Scope and interpretive commitments}
An original property identity is not a scalar index, a linearly ordered value,
or a shorthand for a traditional arithmetic object. In particular,
distinct original occurrences and distinct construction histories remain
distinct even where a declared operation yields the same status.
The symbols in this note are names for distinctions and entailments; they do
not assume classical multiplication or numerical subtraction.

For each named original property $Q$, distinguish witnessed positive evidence
$1_Q$, witnessed negative evidence $0_Q$, unknown evidence $?_Q$,
and conflicting positive/negative reports $\top_Q$. The negative
reference $0_Q$ is \emph{of the original $Q$}, not a second primitive.
Unknown is never silently taken as absence; contradiction is never silently
resolved. Evidence is source- and context-scoped.

The construction includes a recoverable source ledger, declarations of
orientation, versioned joining laws, and witnessed realized relationships.
Possibility or admissibility is not itself a witnessed realization.

\section{Primitive observational domain}
\begin{definition}[Original occurrence]
An occurrence has an original property identity, a value-property identity
where applicable, a source, orientation/role information, observed witnesses,
and provenance. Equality of value-property payloads never by itself merges
distinct occurrence or source identities.
\end{definition}

\begin{definition}[Evidence state]
For a fixed property, source, and observation context, an evidence state is
one of $\{\mathrm{present},\mathrm{absent},\mathrm{unknown},\mathrm{conflict}\}$.
The judgments below use an uncontested witnessed absence only where the
relevant negative report is established and no conflicting positive report
is ignored. Historical reports remain in the ledger even when new reports
change the currently licensed conclusions.
\end{definition}

\begin{definition}[Joining law and obligation]
A joining law $J$ specifies version, scope, port identities and orientations,
source/value admissibility conditions, closure arrangements, witnessed
fulfillment conditions, and any capacity or exclusivity constraints.
An obligation $m$ is a named requirement of an arrangement under $J$,
not an independent operation allocated for the absence of its required
property.
\end{definition}

\begin{definition}[Residual condition]
Relative to an established law and observation snapshot, a residual
condition records what is \emph{not yet fulfilled}: conclusively
defeated, pending, or unresolved through conflicting evidence.
Its recoverable meaning is determined by original occurrences, applicable
laws, witnessed observations and prior realized joins. A residual condition
need not have a separately allocated runtime object.
\end{definition}

\section{CRPP-1: preservation and consequence propagation}
Fix an obligation $m$ belonging to an arrangement $a$ of a closure $M$.
Assume its declared law requires uncontested presence of original property
$Q$ in the relevant context. For a witnessed negative report on that very
requirement, soundness of the requirement yields
\begin{equation}\label{eq:defeat}
  0_Q\land\Req(m,Q)\ \vdash\ \Def(m).
\end{equation}
This does not entail the disappearance of the original surviving occurrences,
nor that every alternative arrangement for $M$ is defeated.

\begin{proposition}[Arrangement-sensitive nonclosure]
Suppose the law enumerates a \emph{complete} family $\mathcal A_M$
of alternative closure arrangements, and each arrangement has sound
declared necessary obligations. Then
\[
 \bigl(\forall a\in\mathcal A_M,\ \Def(a)\bigr)
 \ \Longrightarrow\ \operatorname{Unavailable}(M).
\]
Here $\Def(a)$ requires evidence conclusively defeating at least one
necessary condition of $a$. The converse requires a separate completeness
condition asserting that every way for $M$ to close appears in
$\mathcal A_M$ with the relevant adequacy of its requirements.
\end{proposition}
\begin{proof}
Every admitted closure must realize one arrangement in the complete family.
By hypothesis each such arrangement violates an established necessary
condition. Therefore none of the admitted arrangements can close. This
conclusion is relative to the law's declared scope and evidence snapshot.
\end{proof}

\begin{theorem}[CRPP-1: conditional residual propagation]
Let a finite acyclic network of joining dependencies contain
sound requirement implications, scoped witnesses, complete arrangement
declarations wherever total nonclosure is claimed, and all relevant
source/orientation/capacity conditions. A conclusively defeated
obligation propagates defeat only to dependent arrangements whose
necessary conditions include it. Every downstream nonclosure obtained
by these sound steps is justified by the original witnesses and laws
without construction of an independent absent-counterpart operation
or hypothetical successful intermediate joins.
\end{theorem}
\begin{proof}
At the first dependency, the requirement semantics justify
Eq.~\eqref{eq:defeat}. For an acyclic network, take a topological
ordering. If a predecessor's required arrangement is conclusively
defeated, each dependent arrangement requiring that predecessor is
defeated by its declared necessary condition. A closure is unavailable
only after all its admitted arrangements have been conclusively
defeated. Induction along the ordering preserves soundness.
The deduction uses established conditions, not unrealized output
histories.
\end{proof}

\begin{remark}[No elimination of evidence]
A defeated obligation is not the same as a missing observation,
and an unavailable \emph{particular} closure is not proof that all
other constructions using the surviving occurrences are unavailable.
A change in relevant law version or evidence scope invalidates the
old consequence unless renewed or independently warranted.
\end{remark}

\section{CRPP-2: joining surviving residual distinctions}
Consider original histories $A$ and $B$, surviving occurrences
$a\in A$, $b\in B$, and a further declared law $J_3$.

\begin{definition}[Admissibility and realization]
$\Adm(J_3,a,b)$ means the orientations, source/value types, scope, capacity,
and evidence constraints admit this oriented participation.
$\Real(J_3,a,b,w)$ means a witness $w$ establishes that this relationship
\emph{actually occurred}. Admissibility does not imply realization.
\end{definition}

\begin{definition}[Fulfillment certificate]
A declared law for a prior obligation $m$ specifies when a realized
join witness supplies the required evidence. Denote this by
$\Sat(m\mid J_3,a,b,w)$. It is not inferred merely from similar
property names.
\end{definition}

\begin{proposition}[Scoped obligation fulfillment]
If a realized $J_3(a,b)$ has a witness satisfying the declared
fulfillment condition of a prior obligation $m$, then $m$ is
fulfilled in the resulting construction. Fulfillment of a distinct
obligation $n$ requires its own declared justification.
\end{proposition}
\begin{proof}
This is the soundness condition of $m$'s original fulfillment law.
The realized witness supplies the relevant evidence. No conclusion
about $n$ follows until its own conditions are established.
\end{proof}

\begin{theorem}[CRPP-2: conditional residual joining]
Let $A,B$ have recoverable original ledgers, and suppose $J_3$ is
admissible and \emph{witnessed realized}. Under complete scoped
fulfillment rules, all earlier obligations fulfilled by that witness
may be determined separately from the applicable laws; all remaining
defeated, pending, or conflicting residual conditions are preserved;
and no distinct source history is equated with another. The outcome
does not require independently materializing counterfactual successful
joins for previously unmatched obligations.
\end{theorem}
\begin{proof}
Extend the witnessed ledger with the realized join, participants,
orientation and witness. Each applicable obligation law determines
whether its condition is met by this evidence; other obligations
retain the state justified by their evidence and constraints.
Recoverability of the original ledger distinguishes earlier
nonfulfillment from later fulfillment. All determinations follow
from actual witnesses and laws, not from hypothetical successful
histories.
\end{proof}

\begin{example}[Reciprocal but distinct fulfillment]
$A$ has occurrence $p_A$ and a pending admissible-$Q$ obligation;
$B$ has occurrence $q_B$ and a pending admissible-$P$ obligation.
If the original laws explicitly allow a realized $J_3(p_A,q_B,w)$
to fulfill both, both obligations are satisfied. The sources
and original joining histories remain distinct. If $A$ instead
required a specifically absent $q_A$ of source $S_A$, a later
$q_B$ does not retroactively witness $q_A$'s existence unless
the law permits an alternative fulfillment path.
\end{example}

\section{Information needed for recovery}
Let $\Hist$ be the complete declared original construction ledger.
For a proposed exact representation $E$ and decoder $D$, require
\begin{equation}\label{eq:roundtrip}
       D(E(\Hist))\equiv_{\mathrm{causal}}\Hist.
\end{equation}
The recoverable distinctions include original occurrence/value/source
identities; reports and their contexts; ordering where causally relevant;
law fingerprints and versions; port orientation; witnessed realized joins
and participating occurrences; and source-specific residual obligations.
Derivable statuses may be omitted as \emph{stored objects} if the decoder
recovers them from sufficient retained causes and relations.
Two histories with identical final closure statuses can still differ
causally and must not be merged. For an exact byte-encoding application,
Eq.~\eqref{eq:roundtrip} additionally requires byte-for-byte reconstruction.

\section{Nonconstruction claim and limitations}
The formal claims above establish \emph{logical sufficiency} of existing
law/witness relationships for certain consequences. They do not assert
the absence of all physical transitions or that compiled inference
necessarily outperforms optimized branching, symbolic logic,
incremental evaluation, or a compressed representation.
Instruction count, allocations, memory and energy are experimental
questions distinct from the entailments.

\section{Open derivation: constrained simultaneous joining}
The current results do not determine which of two simultaneously
admissible joins is realized when they compete for the same finite
capacity or change each other's admissibility.
The next proof must define:
\begin{enumerate}[label=(\roman*)]
\item witnessed simultaneity and distinct orientation/causal order;
\item allocation or exclusivity of shared occurrences;
\item composition laws for mutually constraining obligations;
\item necessary and sufficient closure criteria under alternatives;
\item fixed-point semantics for cyclic dependencies that do not assume their own conclusions;
\item exact causal-ledger inversion after joint outcomes.
\end{enumerate}
No unconditional global closure or universal efficiency theorem is
claimed for those unresolved cases.

\section{Experimental contract for a later implementation}
A future constructive evaluator should be independently checked
against a reference evaluator on a fully specified bounded grammar.
Report exact reconstruction, source/witness preservation,
materialized hypothetical states, independently constructed absence
operations, allocations, wall time, memory, and failure cases.
The reference evaluator verifies results; it must not serve as an
unacknowledged source of the constructive law.

\end{document}
